\documentclass{beamer}
\usepackage{hyperref}
\usepackage[T1]{fontenc}
\usepackage{cite}
\usepackage{relsize}
\usepackage{animate}


\usefonttheme[onlymath]{serif}
%C++ symbol
\newcommand{\Rplus}{\protect\hspace{-.1em}\protect\raisebox{.35ex}{\smaller{\smaller\textbf{+}}}}
\newcommand{\Cpp}{\mbox{C\Rplus\Rplus}\hspace{3pt}}
%Scientific notation
\newcommand{\expnumber}[2]{{#1}\mathrm{e}{#2}}

\usepackage{latexsym,amsmath,xcolor,multicol,booktabs,calligra}
\usepackage{graphicx,tikz,listings,stackengine}
\usetikzlibrary{shapes,snakes,matrix, positioning,arrows.meta}
\setbeamertemplate{bibliography item}[text]

\author{Rafael De Luna Loredo}
\title{Programaci\'on en Python y \Cpp}
\subtitle{Funciones}
\institute{Facultad de Ciencias, UASLP}
\date{}
\usepackage{NJFU}

% defs
\def\cmd#1{\texttt{\color{red}\footnotesize $\backslash$#1}}
\def\env#1{\texttt{\color{blue}\footnotesize #1}}
\definecolor{deepblue}{rgb}{0,0,0.5}
\definecolor{deepred}{rgb}{0.6,0,0}
\definecolor{deepgreen}{rgb}{0,0.5,0}
\definecolor{halfgray}{gray}{0.55}

\begin{document}
	
	\AtBeginSubsection[]
	{
		\begin{frame}
			\begin{multicols}{3}
				\tableofcontents[currentsection,subsectionstyle=show/shaded/hide,subsubsectionstyle=show/shaded/hide]
			\end{multicols}
		\end{frame}
	}
	\AtBeginSubsubsection[]{
		\begin{frame}
			\begin{multicols}{3}
				\tableofcontents[currentsection,subsectionstyle=show/shaded/hide,subsubsectionstyle=show/shaded/hide]
			\end{multicols}
		\end{frame}
	}
	
	\AtBeginSection[]
	{
		\begin{frame}
			\begin{multicols*}{3}
				\tableofcontents[currentsection,subsectionstyle=show/shaded/hide,subsubsectionstyle=show/shaded/hide]
			\end{multicols*}
		\end{frame}
	}
	
	\begin{frame}
		\titlepage
		\begin{figure}[htpb]
			\begin{center}
				\includegraphics[width=0.2\linewidth]{fig/logo.png}
			\end{center}
		\end{figure}
	\end{frame}
	
	\begin{frame}{Indice}
		\begin{multicols*}{3}
			\tableofcontents[sectionstyle=show,subsectionstyle=show/shaded/hide,subsubsectionstyle=show/shaded/hide]
		\end{multicols*}
	\end{frame}
	
\section{Introducci\'on}

\begin{frame}{Introducci\'on}
	\begin{itemize}
		\item Es c\'odigo que se encuentra fuera del c\'odigo principal
		\item Nos ayudan a crear c\'odigo mas legible, organizado y f\'acil de mantener
		\item Pueden ser de distintos tipos, correspondientes al tipo de datos
		\item Pueden retornar o no, un valor
		\item Ocupan o no Par\'ametros
		\item Se pueden usar con apuntadores
	\end{itemize}
\end{frame}

\begin{frame}{Sintaxis}
	Una funci\'on se compone de 4 partes
	\begin{itemize}
		\item El tipo de dato de la funci\'on \footnote{En Python no aplica}
		\item El nombre de la funci\'on
		\item Los par\'ametros, si es que lleva
		\item El c\'odigo que ejecuta
	\end{itemize}
\end{frame}

\section{Tipos de Par\'ametros}

\begin{frame}{Por Referencia}
	A la funci\'on se le pasa la direcci\'on de memoria del par\'ametro que solicita, es decir donde se encuentra alojado el valor de dicho par\'ametro por lo que podemos modificar dicho valor dentro de la funci\'on por lo que hay que tener cuidado cuando se use este tipo de par\'ametro ya que podr\'iamos terminar con resultados inesperados
\end{frame}

\begin{frame}{Por valor}
	A la funci\'on se le pasa solo el valor del par\'ametro que solicita, es decir se le pasa una copia de la varaible que se le est\'a pasando, aqu\'i el valor del par\'ametro aunque lo modifiquemos en la funci\'on, en el c\'odigo principal no va a cambiar
\end{frame}

\section{Funciones an\'onimas}

\begin{frame}{Funciones an\'onimas}
	\begin{itemize}
		\item Son funciones sin nombre
		\item Se declaran in situ
		\item Pueden tener como par\'ametros todo el scope
	\end{itemize}
\end{frame}

\section{Recursividad}

\begin{frame}{Recursividad}
	\begin{itemize}
		\item Es m\'as un algoritmo
		\item Aplica para problemas que est\'an definidos sobre si mismos (factorial y fibonacci)
		\item Las funciones se llaman a si mismas
		\item Tienen la desventaja de llegar a consumir un espacio considerable de memoria
		\item Puede ser dif\'icil de implementar
	\end{itemize}
	\centering
	\animategraphics[autoplay,loop,width=4cm]{20}{animaciones/recursividad/recursividad-}{0}{104}
\end{frame}

\section{Macros}
\begin{frame}{Macros}
	\begin{itemize}
		\item Son directivas tratadas por el preprocesador
		\item No son variables ni datos
		\item Sustituyen lo que tengan en donde se haga la llamada a la macro
		\item Disminuyen los tiempos de ejecuci\'on
	\end{itemize}
\end{frame}


\section{Referencias}
\subsection{}
\begin{frame}[allowframebreaks]
	
	\frametitle{Referencias}
	\bibliographystyle{acm}
	\bibliography{ref.bib}
\end{frame}

\end{document}
